% Options for packages loaded elsewhere
\PassOptionsToPackage{unicode}{hyperref}
\PassOptionsToPackage{hyphens}{url}
\documentclass[
  ignorenonframetext,
]{beamer}
\newif\ifbibliography
\usepackage{pgfpages}
\setbeamertemplate{caption}[numbered]
\setbeamertemplate{caption label separator}{: }
\setbeamercolor{caption name}{fg=normal text.fg}
\beamertemplatenavigationsymbolsempty
% remove section numbering
\setbeamertemplate{part page}{
  \centering
  \begin{beamercolorbox}[sep=16pt,center]{part title}
    \usebeamerfont{part title}\insertpart\par
  \end{beamercolorbox}
}
\setbeamertemplate{section page}{
  \centering
  \begin{beamercolorbox}[sep=12pt,center]{section title}
    \usebeamerfont{section title}\insertsection\par
  \end{beamercolorbox}
}
\setbeamertemplate{subsection page}{
  \centering
  \begin{beamercolorbox}[sep=8pt,center]{subsection title}
    \usebeamerfont{subsection title}\insertsubsection\par
  \end{beamercolorbox}
}
% Prevent slide breaks in the middle of a paragraph
\widowpenalties 1 10000
\raggedbottom
\AtBeginPart{
  \frame{\partpage}
}
\AtBeginSection{
  \ifbibliography
  \else
    \frame{\sectionpage}
  \fi
}
\AtBeginSubsection{
  \frame{\subsectionpage}
}
\usepackage{iftex}
\ifPDFTeX
  \usepackage[T1]{fontenc}
  \usepackage[utf8]{inputenc}
  \usepackage{textcomp} % provide euro and other symbols
\else % if luatex or xetex
  \usepackage{unicode-math} % this also loads fontspec
  \defaultfontfeatures{Scale=MatchLowercase}
  \defaultfontfeatures[\rmfamily]{Ligatures=TeX,Scale=1}
\fi
\usepackage{lmodern}
\ifPDFTeX\else
  % xetex/luatex font selection
\fi
% Use upquote if available, for straight quotes in verbatim environments
\IfFileExists{upquote.sty}{\usepackage{upquote}}{}
\IfFileExists{microtype.sty}{% use microtype if available
  \usepackage[]{microtype}
  \UseMicrotypeSet[protrusion]{basicmath} % disable protrusion for tt fonts
}{}
\makeatletter
\@ifundefined{KOMAClassName}{% if non-KOMA class
  \IfFileExists{parskip.sty}{%
    \usepackage{parskip}
  }{% else
    \setlength{\parindent}{0pt}
    \setlength{\parskip}{6pt plus 2pt minus 1pt}}
}{% if KOMA class
  \KOMAoptions{parskip=half}}
\makeatother
\usepackage{longtable,booktabs,array}
\usepackage{caption}
\captionsetup[table]{skip=6pt}
\newcounter{none} % for unnumbered tables
\usepackage{calc} % for calculating minipage widths
\usepackage{caption}
% Make caption package work with longtable
\makeatletter
\def\fnum@table{\tablename~\thetable}
\makeatother
\setlength{\emergencystretch}{3em} % prevent overfull lines
\providecommand{\tightlist}{%
  \setlength{\itemsep}{0pt}\setlength{\parskip}{0pt}}
% Auto-generated by infrastructure/rendering/_pdf_latex_helpers.py::extract_math_font_preamble.
% Minimal header passed to Pandoc via -H so Beamer slides pick up the same math font
% as the combined-PDF path. Adjust the manuscript's preamble.md to change the font globally.
\usepackage{unicode-math}
\setmathfont{latinmodern-math.otf}

% Manuscript macro/package fallbacks (newcommand -> providecommand).
% Core mathematics
\usepackage{amsmath}
\usepackage{amssymb}
% Document layout
% Tables
\usepackage{booktabs}
% Algorithm typesetting (for pseudocode in §2 Methodology)
% Code listings
\usepackage{listings}
% Typography and formatting
% Cross-references and citations
% ── Unicode-capable mono font for code listings ──────────────────────
% JuliaMono (TeX Live 2026) covers the full math/Greek glyph set used in
% scientific code blocks; the default lmmono lacks \alpha/\mu/\partial/\nabla.
%
% The hardcoded TeX Live 2026 path below is portable to most macOS / Linux
% TeX Live installs that placed JuliaMono under the default texmf-dist path.
% If your TeX Live version differs (e.g., 2025, 2027) or your distribution
% installs fonts elsewhere, comment out the explicit Path = line — fontspec
% will then search the system font path. If JuliaMono is not installed at
% all, comment out the whole block; xelatex falls back to lmmono with
% reduced Unicode coverage but the document still compiles.
% Math font for unicode-math: Latin Modern Math (TeX Live) has full BMP
% coverage including U+2223 (\mid), U+226A/226B (\ll/\gg), and the Greek/
% blackboard letters used in equations. Without an explicit \setmathfont,
% unicode-math falls back to lmroman text font which lacks several glyphs
% and emits "Missing character" warnings on every \mid in math mode.

% Slide layout defaults for warning-clean scientific decks.
\usepackage{etoolbox}
\IfFileExists{xurl.sty}{\usepackage{xurl}}{}
\IfFileExists{seqsplit.sty}{\usepackage{seqsplit}}{\newcommand{\seqsplit}[1]{#1}}
\protected\def\breakseq#1{\seqsplit{#1}}
\protected\def\breaktt#1{\begingroup\ttfamily\seqsplit{#1}\endgroup}
\setlength{\emergencystretch}{6em}
\tolerance=5000
\hbadness=10000
\hfuzz=1pt
\setlength{\tabcolsep}{2pt}
\AtBeginEnvironment{longtable}{\tiny\renewcommand{\arraystretch}{0.86}\setlength{\tabcolsep}{1pt}}
\AtBeginEnvironment{tabular}{\tiny\renewcommand{\arraystretch}{0.86}\setlength{\tabcolsep}{1pt}}
\AtBeginEnvironment{equation}{\tiny}
\AtBeginEnvironment{equation*}{\tiny}
\AtBeginEnvironment{align}{\tiny}
\AtBeginEnvironment{align*}{\tiny}
\AtBeginEnvironment{itemize}{\footnotesize}
\AtBeginEnvironment{enumerate}{\footnotesize}
\AtBeginEnvironment{description}{\footnotesize}
\setbeamerfont{caption}{size=\tiny}
\setbeamerfont{caption name}{size=\tiny}
\setbeamerfont{normal text}{size=\small}
\setbeamerfont{frametitle}{size=\small}
\setbeamerfont{section title}{size=\footnotesize}
\setbeamerfont{subsection title}{size=\footnotesize}
\setbeamertemplate{section page}{%
  \centering
  \begin{beamercolorbox}[sep=12pt,center,wd=\paperwidth]{section title}
    \parbox{0.86\paperwidth}{\centering\usebeamerfont{section title}\insertsection\par}
  \end{beamercolorbox}
}
\setbeamertemplate{subsection page}{%
  \centering
  \begin{beamercolorbox}[sep=8pt,center,wd=\paperwidth]{subsection title}
    \parbox{0.86\paperwidth}{\centering\usebeamerfont{subsection title}\insertsubsection\par}
  \end{beamercolorbox}
}
\setlength{\abovecaptionskip}{2pt}
\setlength{\belowcaptionskip}{0pt}

% Natbib and cross-reference fallbacks — slides don't load natbib
% or cleveref, but combined-PDF manuscript prose may emit these
% commands. The fallback renders citations as a bracketed key list
% and cross-references as detokenized labels so slides don't fail on
% undefined control sequences or raw underscores. \providecommand is
% a no-op if packages load later.
\providecommand{\citep}[1]{[#1]}
\providecommand{\citet}[1]{#1}
\providecommand{\citealp}[1]{#1}
\providecommand{\citeauthor}[1]{#1}
\providecommand{\citeyear}[1]{#1}
\providecommand{\cref}[1]{\texttt{\detokenize{#1}}}
\providecommand{\Cref}[1]{\texttt{\detokenize{#1}}}
\providecommand{\crefrange}[2]{\texttt{\detokenize{#1}}--\texttt{\detokenize{#2}}}
\providecommand{\Crefrange}[2]{\texttt{\detokenize{#1}}--\texttt{\detokenize{#2}}}
\providecommand{\paragraph}[1]{\textbf{#1}\ }

\newtheorem{proposition}{Proposition}
\newtheorem{hypothesis}{Hypothesis}
\newtheorem{remark}[theorem]{Remark}
\newtheorem{axiom}{Axiom}
\newtheorem{property}{Property}
\usepackage{bookmark}
\IfFileExists{xurl.sty}{\usepackage{xurl}}{} % add URL line breaks if available
\urlstyle{same}
% fallback for those not using the hyperref driver hyperxmp:
\makeatletter
\@ifundefined{xmpquote}{\newcommand{\xmpquote}[1]{#1}}{}
\makeatother
\hypersetup{
  hidelinks,
  pdfcreator={LaTeX via pandoc}}

\author{\texorpdfstring{}{}}
\date{}

\begin{document}

\section{Operator OpSec: Identity, Egress, and Incident Response for
Agent Work}\label{sec:opsec}

\begin{frame}[allowframebreaks]{Operator OpSec: Identity, Egress, and
Incident Response for Agent Work}
Operating-system compartmentalization constrains what a compromised
component can reach. Operator operational security constrains what a
correctly functioning agent is allowed to do. The threat model of
sec.~3 separates exploitation from authorized
misuse, and the two defenses divide along that line: isolation
architectures raise the cost of the exploitation path, while OpSec is
the operator-side discipline for the authorized-misuse path, where no
kernel exploit is necessary because the agent already holds a valid
credential, an enabled integration, and a standing authorization. The
AISI incident report makes the distinction concrete in its own incident
report: agents took out-of-scope internet actions in 10 of 122 test runs
without escaping the sandbox, because internet access was intentionally
available and the authority to act had already been granted (Institute
2026). This section treats OpSec as the reviewable extension of the
\end{frame}

\begin{frame}[allowframebreaks]{Operator OpSec: Identity, Egress, and
Incident Response for Agent Work}
trust-domain architecture of sec.~10: the 7
domains and 9 controls define what the system separates; operator
discipline determines whether those separations survive contact with
daily work. Three developments sharpen that discipline into something
more than exhortation: standards work now defines what an agent's
identity can look like, incident reporting documents credential theft as
a recurring failure mode, and transparency regulation turns attribution
hygiene into a compliance surface.

\end{frame}

\begin{frame}[allowframebreaks]{Operator OpSec: Identity, Egress, and
Incident Response for Agent Work}
Structured self-assessment gives that discipline a reviewable form:
operator posture assessments and human-oversight checklists, published
as part of a practical deployment guide for cognitive-integrity
controls, turn the review of agent-facing authority into a repeatable
procedure rather than an ad hoc audit (Friedman 2026). The subsections
below follow the same shape --- each practice is stated as a checkable
configuration or a rehearsed response, not as advice.
\end{frame}

\begin{frame}[fragile,allowframebreaks]{Identity compartmentalization
across agent contexts}
\protect\phantomsection\label{identity-compartmentalization-across-agent-contexts}
The trust-domain design separates \breaktt{personal\_identity} from
\texttt{agent\_execution} (tbl.~8). That
separation exists at the operating-system layer only if it also exists
at the identity layer. An operator who signs in to an agent-driven
workflow with a personal single-sign-on identity has fused the two
domains regardless of how well the machine is compartmentalized: the
agent's ambient session then speaks for the person, and every boundary
inside the machine becomes an implementation detail the credential route
ignores.

\end{frame}

\begin{frame}[fragile,allowframebreaks]{Identity compartmentalization
across agent contexts}
Four practices keep the domains distinct. First, dedicated accounts for
agent-mediated services, enrolled in their own multi-factor factors, so
that agent work never borrows the operator's personal identity. Second,
no ambient production session inside an agent execution environment:
cached browser sessions, cloud CLI tokens, and ssh-agent keys stay in
the domains that own them. Third, identity lifetime aligned with task
lifetime: an account minted for a task expires with the task, which caps
the value of any exfiltrated session material. Fourth, deliberate
handling of the human-authorized transfer paths that compartmentalized
systems deliberately provide --- clipboard export, cross-domain copy,
shared browser profiles --- because these are user-approved crossings of
the very boundary the identity plan maintains (Q. O. Project 2026b). An
attacker who persuades the operator to move an authenticated session
\end{frame}

\begin{frame}[fragile,allowframebreaks]{Identity compartmentalization
across agent contexts}
into the agent's domain has attacked the workflow, not escaped the
hypervisor, and no operating-system choice reverses that.
\end{frame}

\begin{frame}[allowframebreaks]{Agent identity: workload identity, token
exchange, and downscoped tokens}
\protect\phantomsection\label{agent-identity-workload-identity-token-exchange-and-downscoped-tokens}
The identity practices above are currently operator discipline because
no ratified standard yet defines how an autonomous agent authenticates
as itself. Standards work is converging on an answer. An IETF draft for
AI-agent authentication composes three existing mechanisms: workload
identity from the SPIFFE and WIMSE families, so an agent holds a
cryptographically attestable workload identity rather than a borrowed
user account; OAuth 2.0 token exchange (RFC 8693), so a task receives
its own credential by exchanging an attestable identity against policy;
and transaction tokens, so each credential binds a specific subject,
audience, and set of requested authorities for a limited time (Force
2026). The composition is not yet a ratified standard, and this review
does not forecast its ratification. Its shape, however, is already the
\end{frame}

\begin{frame}[allowframebreaks]{Agent identity: workload identity, token
exchange, and downscoped tokens}
shape the trust-domain design asks for: an agent identity that is
distinct from the operator's, minted per task, bound to a specific
audience, narrowed to the authorities the task requires, and
short-lived.

Short-lived downscoped tokens are the operational core of that shape,
and they can be built today without waiting for the draft: a
token-issuing service exchanges a workload attestation for a credential
whose scope is the task's tool grants and whose expiry is the task's
lifetime. This is the scoped-credentials control of
tbl.~9 implemented at the identity layer, and it
composes with the orchestration-layer capability handoffs of
sec.~13: the token an executor presents to a tool
broker is exactly the narrow, unforgeable reference the delegation
analysis there demands. Until such machinery is standard, the fallback
remains the four practices above --- and the design intent is worth
recording now, because operators who structure agent work around
\end{frame}

\begin{frame}[allowframebreaks]{Agent identity: workload identity, token
exchange, and downscoped tokens}
distinct, downscoped, task-lifetime identities will need no
re-architecture when the standards land.
\end{frame}

\begin{frame}[allowframebreaks]{Data-loss path analysis}
\protect\phantomsection\label{data-loss-path-analysis}
Every egress-capable destination is a disclosure decision, and the set
is larger than it looks. Model providers receive the task context
itself: prompts, attached documents, repository contents, error
messages. Repository forges receive pushed branches and issue text. CI
systems receive logs that often embed environment variables and secrets.
Error-reporting and telemetry channels receive exception payloads.
Messaging integrations receive whatever an agent decides to summarize. A
hostname allowlist is not a complete data-loss policy, because a
permitted destination that itself accepts uploads or messages is a data
destination in its own right (sec.~10).

\end{frame}

\begin{frame}[allowframebreaks]{Data-loss path analysis}
The sharpest instance is the model provider itself. The decision to
transmit private code or documents to a hosted model is a separate
trust-boundary decision that no operating-system choice can resolve:
once transmitted, retention policy, staff access, training use, and the
provider's own breach exposure are outside the workstation's authority
entirely. This is a decision to classify material before it enters a
task environment, not after; after an agent has read private material,
assume it can leave by any egress the task's tools permit. Contractual
terms about data handling are policy commitments, not containment. Local
inference changes the transport but not the structure of the decision:
the model service, the agent execution domain, and credential mediation
should remain separated with explicit authentication between them,
because shared compute is not a reason to share trust.

\end{frame}

\begin{frame}[allowframebreaks]{Data-loss path analysis}
Official guidance now frames the same decision from the adoption side.
The CISA-led Five Eyes guidance on careful adoption of agentic AI
services recommends sandboxed deployment, threat-model-based evaluation,
red-teaming, and --- most directly relevant to the egress decision ---
limiting early agent deployments to low-risk tasks until the surrounding
controls are demonstrated (Cybersecurity and Agency 2026). Low-risk-task
staging is the operational complement to the data-loss analysis above: a
task that cannot reach private material cannot disclose it, so the
material classification and the task's risk tier should be decided
together, before the first agent run, not reconstructed after an
incident.
\end{frame}

\begin{frame}[allowframebreaks]{Tool-bridge grants as the OpSec surface}
\protect\phantomsection\label{tool-bridge-grants-as-the-opsec-surface}
Tool bridges --- the filesystem, browser, repository, CI, cloud, and
messaging integrations through which agents act --- are the OpSec
surface proper. Each bridge is a standing grant of authority: a
filesystem root the agent may read or write, a browser profile it may
drive, a repository it may push to, a CI pipeline it may trigger, a
cloud API scope it may exercise, a messaging channel it may post to.
These grants must be reviewed as grants, not as features. An isolated
process with a powerful API token is still a powerful actor:
disposability of the environment does not dispose of the grant, and an
agent inside a disposable qube holding an owner-scoped cloud credential
carries the owner's cloud authority until that credential expires or is
revoked (Q. O. Project 2026a).

\end{frame}

\begin{frame}[allowframebreaks]{Tool-bridge grants as the OpSec surface}
Complete mediation (Saltzer and Schroeder 1975) is the governing
principle: every privileged operation crosses a checked channel, and the
tool broker of sec.~13 is the natural enforcement
point. Operator discipline makes the grants enumerable, task-scoped,
revocable, and logged. The review question for each bridge is the
authority question of sec.~4 asked one
layer up: what can this integration read, transmit, sign, or change, and
does any single task actually need all of that?
\end{frame}

\begin{frame}[allowframebreaks]{Attribution and footprint hygiene}
\protect\phantomsection\label{attribution-and-footprint-hygiene}
Agent traffic is identifiable. Model API endpoints, client user agents,
commit metadata, issue formatting, and activity timing all distinguish
agent-driven work from ordinary operator work. Footprint hygiene means
deciding deliberately what agent activity discloses, rather than letting
defaults decide. Use separate accounts or tenants for agent-driven
commits, issues, and messages, so that a compromised or mistaken agent
cannot speak with the operator's personal voice or reach the audiences
tied to the personal identity. Keep provenance accurate rather than
laundered: an agent-generated change should be recorded as
agent-generated, because obscuring agent authorship destroys exactly the
record that incident response and the audit obligations of
sec.~13 depend on. This hygiene is credential
protection, not anonymity: the goal is that agent activity cannot be
\end{frame}

\begin{frame}[allowframebreaks]{Attribution and footprint hygiene}
replayed as the operator's identity --- a different requirement from
concealing that the activity occurs, which is the distinct province of
the anonymity-focused systems reviewed in sec.~9.

Attribution hygiene is also becoming a regulated surface. The European
Union's AI Act transparency obligations, applicable from 2 August 2026,
require that people interacting with AI systems be informed of the AI
nature of the interaction and --- where autonomous agents act on
someone's behalf --- of the natural person or legal entity on whose
behalf the agent acts; the European Commission's guidelines name
AI-agent disclosure duties explicitly (Commission 2026). An operator
deploying agents into EU-reachable workflows therefore needs the
attribution layer to answer two questions by design: that an AI system
is acting, and whose authority it acts under. This coincides with,
rather than conflicts with, the hygiene above: the accurate provenance
\end{frame}

\begin{frame}[allowframebreaks]{Attribution and footprint hygiene}
record that incident response needs is the same record that disclosure
compliance requires, and both are served by making agent identity
explicit at the identity layer rather than reconstructing it from logs.
Operators outside the regulation's reach still face the same design
requirement wherever counterparties, auditors, or platforms demand
equivalent disclosure.
\end{frame}

\begin{frame}[fragile,allowframebreaks]{Credential hygiene: short-lived,
service-scoped}
\protect\phantomsection\label{credential-hygiene-short-lived-service-scoped}
The credential plan follows the scoped-credentials control: prefer
short-lived, repository- or service-specific credentials issued per task
and expiring with it, and never hand an agent the owner's general cloud
identity merely because a task sometimes needs a deployment
(tbl.~9). The practical compromise is ambient session
tokens. Cached cloud CLI credentials, browser session cookies, and
long-lived personal access tokens already exist in the environment, are
easy to forget, and are precisely what an agent picks up when it
inspects its surroundings. The defense is structural rather than
exhortative: keep ambient tokens out of the \texttt{agent\_execution}
domain instead of trying to teach the agent restraint. Declarative
environments add their own trap on this point: the Nix store is readable
by all users, and secrets embedded in configuration can reach external
\end{frame}

\begin{frame}[fragile,allowframebreaks]{Credential hygiene: short-lived,
service-scoped}
caches, so the documented practice is to read secrets at runtime under
access control (N. Project 2026). Configuration cleanliness must not
come at the cost of credential exposure.
\end{frame}

\begin{frame}[allowframebreaks]{Credential theft as the observed failure
mode}
\protect\phantomsection\label{credential-theft-as-the-observed-failure-mode}
The credential-theft pattern is no longer hypothetical. Anthropic's
September 2026 threat-intelligence report --- the vendor's account of
activity observed on its own platform --- documents threat actors
obtaining legitimate API keys by theft and credential abuse, then
driving coding and agent tooling under the stolen identity; the report's
designated threat-activity identifiers GTG-50020 and GTG-50021 describe
key-harvesting operations against agent platforms, and the report
assesses that ``vibe hacking'' --- attackers using coding agents
directly in extortion and intrusion work --- persists as the top
observed AI-enabled threat (Anthropic 2026b). The pattern matters to
operator OpSec because it is the authorized-misuse path executed with
the operator's own authority: an attacker holding a valid agent-platform
\end{frame}

\begin{frame}[allowframebreaks]{Credential theft as the observed failure
mode}
API key needs no exploit at all, only the standing grants that key
carries. Every credential control in this section exists to shrink
exactly that surface --- short-lived task-scoped issuance so a stolen
key's value decays in hours, dedicated accounts so a stolen agent
credential cannot reach personal identity, per-key scope review so a
stolen key cannot reach more than its task, and egress monitoring on
agent-platform usage so an anomalous consumer of the operator's quota is
visible. The account that pays for agent capability is itself a
high-value credential, and it belongs in the same reviewed inventory as
every tool-bridge grant.
\end{frame}

\begin{frame}[allowframebreaks]{The intake discipline}
\protect\phantomsection\label{the-intake-discipline}
Treat returned code, documents, and build artifacts as untrusted until
checked. A diff produced by an agent is untrusted input into the trusted
codebase; a document it produced is untrusted input into the document
pipeline; a binary it built is untrusted until its construction is
accounted for. No automatic promotion of an agent's home directory,
working tree, or container image into a trusted environment
(tbl.~9, minimal shared state). The vendor evidence
cuts both ways here: DARPA's AIxCC results demonstrate that
machine-generated patches for real vulnerabilities are achievable (DARPA
2025), while the Anthropic campaign investigation --- the investigating
vendor's report --- found agents overstating findings and fabricating
unusable credentials (Anthropic 2026a). Generation quality and
acceptance authority remain separate questions; the review that admits
\end{frame}

\begin{frame}[allowframebreaks]{The intake discipline}
generated material is the same independent review that
sec.~14 requires for generated
policy.
\end{frame}

\begin{frame}[fragile,allowframebreaks]{Rehearsed incident response for
agent-era compromise}
\protect\phantomsection\label{rehearsed-incident-response-for-agent-era-compromise}
Authorized misuse leaves no exploit artifacts: no crash, no dropped
binary, no anomalous fault trace. The response therefore cannot depend
on detecting the compromise before acting, and the working assumption
must be that everything the task could touch was exfiltrated or altered.
The rehearsed sequence is: revoke every credential the task environment
could reach --- a superset of those it demonstrably used; destroy and
rebuild the environment from known-good configuration; treat data
already transmitted as disclosed; restore from independent backups.
Distinguish deleting an environment from revoking the credentials it
could access: deletion without revocation leaves live authority in
unknown hands. Environment refresh between tasks and prompt patching of
management and transfer paths belong to the same discipline --- QSB-118
\end{frame}

\begin{frame}[fragile,allowframebreaks]{Rehearsed incident response for
agent-era compromise}
documented a route from an already-compromised qube to dom0 command
injection under specified conditions, fixed in
\breaktt{qubes-core-dom0-linux} 4.3.22 (Q. O. Project 2026c), and
recovery speed under such advisories is an operator OpSec property, not
only a vendor property. Recovery that has never been rehearsed
(tbl.~9, rehearsed recovery) is a plan-shaped hope. The
credential-theft pattern of the previous subsection sharpens the
response plan in one specific way: revocation must cover the
agent-platform accounts themselves, not only the task-scoped credentials
they issued, because a stolen platform key can mint new task authority
after the original task environment is gone.

\end{frame}

\begin{frame}[fragile,allowframebreaks]{Rehearsed incident response for
agent-era compromise}
Incident-response playbooks published for cognitive-integrity
deployments ground that rehearsal: the revocation-first,
rebuild-from-known-good, treat-transmitted-as-disclosed sequence appears
there as practiced steps with assigned roles and timing rather than as
principles to be improvised under pressure (Friedman 2026).
\end{frame}

\begin{frame}[allowframebreaks]{Zero trust applied to agent tool calls}
\protect\phantomsection\label{zero-trust-applied-to-agent-tool-calls}
Zero-trust architecture removes implicit trust derived from network
location or prior approval: every access request is evaluated per
request against policy, and the policy decision point is separate from
the resource it guards (Rose et al. 2020). Agent tool calls have exactly
this shape. A task's prior approval is not a standing authorization; a
request's origin inside a trusted agent is not evidence of legitimacy;
each call is judged on the operation requested, the artifact named, the
destination addressed, the scope claimed, and the expiry asserted.
Operation mediation is zero trust stated operationally
(tbl.~9): a service that permits ``sign this exact
artifact for this release'' is checkable per request, while ``use the
signing key responsibly'' is not checkable at all. Zero trust for agent
work therefore lands on the same design point as the orchestration
\end{frame}

\begin{frame}[allowframebreaks]{Zero trust applied to agent tool calls}
boundaries of sec.~13 and the configuration
separation of sec.~14: place the
decision point outside the agent, and make the decision context rich
enough to audit after the fact. The agent-identity standards described
above are the direction in which this posture standardizes: workload
identity, token exchange, and transaction tokens are zero trust rendered
as credential machinery, so that the per-request judgment has a
per-request credential to judge.
\end{frame}

\end{document}
